\documentclass{beamer}
\usepackage{verbatim}
\usepackage{wrapfig}
\usepackage{color} 
\usetheme{Amsterdam}
\title{Working with HAIKU for Google Code In 2014}
\author{\texorpdfstring{Alex-P. Natsios\newline\url{drakevr@2f30.org}}{Author}}
\institute{High School of Nikaia, Larissa, Greece}
\date{16 Jan 2015}
\begin{document}
    \begin{frame}
       \titlepage
    \end{frame}

    \begin{frame}{}
        \begin{figure}{}{}
            \includegraphics[scale=0.2]{images/GCI_logo.eps}
        \end{figure}
    \end{frame}

    \begin{frame}{What is Google Code In?}

        {\bf Google Code In}, also called {\bf GCI}, is a Google sponsored contest,
        that enables students to work in real life conditions/workflows with Free
        And Open Source software projects.

        \begin{block}{Important points:}
            \begin{itemize}
                \item GCI happens {\bf Online}
                \item Any Student between the ages of 13 and 17 may enter
                \item Gentle introduction to Open Source development
                \item Open Source organizations create tasks for students to work on
                \item Students are free to choose just those tasks that interest them
                \item Students earn certificates, t-shirts, hooded
                      sweatshirts and can work towards the Grand Prize
                      trip to Google campus
            \end{itemize}
        \end{block}
    \end{frame}

    \begin{frame}{Why should you participate in GCI}
        \begin{itemize}
            \item Apply classroom skills in real life projects
            \item Learn about the many exciting projects you can choose from
            \item Become a part of the community
            \item Easy to enter a project, each organization provides mentors
                  to guide students.
            \item Working with Open Source does not imply coding, there are
                  many other equally important tasks to choose from
        \end{itemize}
    \end{frame}

    \begin{frame}{Types of Tasks}
        Each task is designed to take approximately the same amount of time to
        complete (about 3-5 hours) 

        \begin{block}{Categories of tasks}
            \begin{itemize}
                \item Coding
                \item Documentation/Training
                \item Outreach/Research
                \item Quality Assurance (QA)
                \item User Interface (UI)
            \end{itemize}
        \end{block}
    \end{frame}

    \begin{frame}{}
        \begin{figure}{}{}
            \includegraphics[scale=0.5]{images/HAIKU_logo_transparent_black.eps}
        \end{figure}
    \end{frame}

    \begin{frame}{What is HAIKU}
        \begin{itemize}
            \item HAIKU is a fast, efficient, easy to use and learn
                  open source operating system that specifically targets 
                  Personal Computing.
            \item Custom kernel designed for responsiveness
            \item Unified, cohesive integrated interface
            \item Best defaults, not maximum configurability
            \item Open Source (MIT licensed)
        \end{itemize}
    \end{frame}

    \begin{frame}{Goals}
        \begin{itemize}
            \item Simplicity is the key
            \item Multiprocessors
            \item Multimedia (Exciting and HIGH QUALITY graphics and audio stacks)
            \item Desktop Focus
            \item small FAST system and startup
        \end{itemize}
    \end{frame}

    \begin{frame}{What does that mean for me? How will HAIKU help me?}
        \begin{itemize}
            \item A mature project in which fresh ideas are continuously tested
            \item A develpment playground to know about new stuff and aquire
                  much needed skills in todays industry
            \item A friendly place to learn about how computers and Operating Systems work
            \item Even a small contribution goes a long way and shows value both
                  for the project and yourself (especialy if you plan to be a
                  programmer in the future).
        \end{itemize}
    \end{frame}

    \begin{frame}{Cool! How can I help?}
        \begin{itemize}
            \item Try it out, report bugs
            \item Lots of bugs to fix, mostly C/C++
            \item Translate into other languages
            \item Help out with porting open source apps at HaikuPorts
            \item Write new applications
            \item Spread the word
        \end{itemize}
    \end{frame}

    \begin{frame}{Links}
        HAIKU homepage:\\
        \ \ \url{http://haiku-os.org}\\
        Software for HAIKU:\\
        \ \ \url{http://haikuports.org}\\
        Facebook:\\
        \ \ \url{https://www.facebook.com/pages/Haiku-OS}\\
    \end{frame}

    \begin{frame}{Thank you for your attention!}
        \center Alex-P. Natsios\\
        \center\url{drakevr@2f30.org}\\
        \center\url{http://drakevr.gr}\\
        \center\url{http://www.linkedin.com/in/drakevr}\\
        \center\url{http://www.github.com/drakevr}\\
        \center\url{http://www.facebook.com/drakevr}\\
        \center\url{http://www.twitter.com/drakevr}
    \end{frame}

\end{document}
